This laboratory connected sensor readings with decisions and actuator responses
in eight simulated or physical systems. The experiences show that a circuit
diagram alone does not establish correct behavior: thresholds, timing,
priorities, electrical limits and the measurements used to check them also
matter. The following conclusions address each system and its next check.

The LED-bar experience showed how a potentiometer reading can be converted
from the ADC range into a percentage and a discrete number of illuminated
LEDs. Because the display has only ten steps, it is recommended to state the
rounding rule and compare the serial value with the LED count at the minimum,
midpoint and maximum input.

The four-LDR system combined several light readings while treating a failed
channel and loss of camera supply as separate alarm conditions. A mean is
useful only when invalid channels are excluded and at least one valid reading
remains. It is recommended to document the divider orientation and light
threshold, then test each fault independently from ordinary darkness.

The occupancy system used PIR events and an exit button to estimate a count,
while the gas alarm took priority over the normal LCD message. A missed or
repeated event can make that estimate differ from the real occupancy. It is
recommended to check the five-second lockout, the count limits and the gas
alarm both separately and when the room is full.

The multiparameter station showed how distance, temperature and humidity can
share a two-line LCD through alternating views. The warning must remain
legible whichever view is active. It is recommended to check the three-second
switching interval and trigger the distance and temperature alerts separately.

The motor simulations showed markedly different speed readings for the three
connections reported in the development section. These simulation values do
not establish that an Arduino output can safely power a physical motor. It is
recommended to verify the instrument units in the screenshots and use a
suitable driver and external supply in a physical circuit.

The LED-dimming experience distinguished gradual PWM control from the on/off
response of a pin without PWM support. It is recommended to compare the ADC
reading, PWM command and visible output at the same potentiometer settings,
while retaining a current-limiting resistor in the LED circuit.

The relay experience separated the motor supply path from the Arduino control
signal, but the active input level depends on the particular module. It is
recommended to verify the COM/NO contacts, supply voltage, active level and
motor protection before using the demonstration as evidence of correct
switching.

The irrigation system combined calibrated moisture thresholds with a timed
pump cycle and a mandatory idle interval. Its decisions depend on the measured
dry and wet endpoints, so it is recommended to record those readings and check
the dry, normal and excessive-moisture states as well as the 10-s and 15-s
transitions. Water savings cannot be quantified without measuring consumption.

Overall, the laboratory illustrates why simulation, physical observation and
recorded measurements should support one another. Future iterations should
retain the actual thresholds, calibration values and timing observations with
each checkpoint so that the reported behavior can be reproduced and checked.
